\section{March 10: divisors}

A goal for us in the coming weeks is to define intersection products for Chow groups.

\begin{example}
  Let $X$ be a topological space. 
  The cap product is defined on singular homology and cohomology on $X$ as follows:
  \[ H_i (X,\bbZ) \ten H^j (X,\bbZ) \to H_{i-j}(X,\bbZ), \quad (\sigma, \psi) \mapsto \psi(\sigma). \]
  The cohomology group $H^j (X,\bbZ)$ may be regarded as the collection of operations on $H_\bullet (X, \bbZ)$ with grading shift $-j$.
\end{example}

For a scheme $X$, we will define 
\[ \CH_i (X) \ten  \text{``}\CH^j (X)\text{''} \to \CH_{i-j} (X), \]
where ``$\CH^j (X)$'' contains the $j$-th Chern class of vector bundles on $X$.
The simplest case is $j = 1$, corresponding to line bundles or Cartier divisors.

\subsection{Divisors}

  \begin{definition}
    Let $X$ be a $n$-dimensional variety.
    A \emph{Weil divisor} on $X$ will mean an $(n-1)$-cycle on $X$.
    A \emph{Cartier divisor} on $X$ is defined by data ${(U_\alpha, f_\alpha)}$ where ${U_\alpha}$ is an open cover of $X$ and $f_\alpha$ is a non-zero rational function on $U_\alpha$ such that $f_\alpha / f_\beta$ is a unit on $U_\alpha \cap U_\beta$.
    In other words, it is an element of 
    \[ \Div (X) \coloneqq \Gamma (X, \fkk(X)^* / \mathcal{O}_X^*). \]
    A Cartier divisor is called \emph{principal} if it lies in the image of the natural map 
    \[ \Gamma (X, \fkk(X)^*) \to \Div (X) . \]
    Two Cartier divisors are called \emph{linearly equivalent} if their difference is principal, and we write $D \sim D'$.
  \end{definition}

  Note 
  \[ 0 \to \mathcal{O}_X^* \to \fkk(X)^* \to \Div (X) \to 0 \] is exact, hence we have the induced map
  \[ \Div (X) \to \Pic (X) \cong H^1 (X, \mathcal{O}_X^*). \]
  We have the natural map $\Div (X) \to Z_{n-1} (X)$ sending a Cartier divisor $\{(U_\alpha, f_\alpha)\}$ to a Weil divisor $\sum_V \ord_V (f_\alpha) \cdot V$, where $V$ runs through all irreducible subvarieties of codimension 1 in $X$ and $\ord_V (f_\alpha)$ is the order of vanishing of $f_\alpha$ along $V$ whenever $V \cap U_\alpha \neq \emptyset$.
  Since $f_\alpha / f_\beta$ is a unit on $U_\alpha \cap U_\beta$, $\ord_V (f_\alpha) = \ord_V (f_\beta)$, so the above definition is independent of the choice of $\alpha$.

  \begin{example}
    If $X$ is a smooth variety, $\Pic (X) \cong \CH_{n-1} (X)$ (see \cite[Chapter II, Section 6]{Har77}).
    In general, the map $\Pic (X) \to \CH_{n-1} (X)$ is neither injective nor surjective.
    For instance, $X \coloneqq \{x y = z^2\} \subset \bbA^3$ and $l \coloneqq \{x = z = 0\} \subset X$ is a Weil divisor that is not Cartier.
    And $\CH_1 (X) = \Cl (X) = \bbZ/2$ generated by $[l]$.
    This implies $\Pic(X) \to \CH_1 (X)$ is the zero map.
    In fact, this is injective (using $X$ is normal), so $\Pic(X) = 0$.
  \end{example}

  Note that the definition of Cartier divisors extends to schemes.
  (But $\left(\Div(X)/\sim\right) \to \Pic (X)$ is NOT an isomorphism in general!)

\subsection{Intersection with divisors}

  Line bundles always can be pull-backed, so do Cartier divisors modulo linear equivalence.

  \begin{example}
    Let $\pi: \widetilde{\bbA^n} \to \bbA^n$ be the blow-up of $\bbA^n$ at the origin, and let $\bbP^{n-1} \cong E \subset \widetilde{\bbA^n}$ be the exceptional divisor.
    Then $\calO_{\bbA^n}(E)|_E = \calN_{E/\widetilde{\bbA^n}} = \calO_E(-1)$, but $E$ does NOT pull-back to $E$.
  \end{example}

  For a Cartier divisor $D$ on a scheme $X$, define the \emph{support} of $D$ to be
  \[ |D| \coloneqq \Supp (D) \coloneqq \bigcup_{V \subset X, \ord_V(D) \neq 0} V . \]
  Cartier divisors have supports , but line bundles do NOT.

  \begin{definition}
    A \emph{pseudo-divisor} on $X$ is a triple $(L,Z,s)$ where 
    \begin{itemize}
      \item $L$ is a line bundle on $X$,
      \item $Z \subset X$ is a closed subset of $X$,
      \item $s$ is a nowhere vanishing section on $L|_{X-Z}$.
    \end{itemize}
  \end{definition}

  A pseudo-divisor is either
  \begin{itemize}
    \item $(L,X,1)$ for some line bundle $L$ on $X$, or
    \item $(\calO_X (D), |D|, s_D)$ for some Cartier divisor $D$ on $X$, where $s_D$ corresponds to $1 \in \fkk(X)^*$ with $\calO_X (D) \subset \fkk(X)^*$.
  \end{itemize}
  We also denote a pseudo-divisor $D$ and the corresponding triple by $(L_D, |D|, s_D)$.

  \begin{definition}
    For pseudo-divisors $D$ and $D'$, define
    \begin{align*}
      D + D' &\coloneqq (L_D \ten L_D', |D| \cup |D'|, s_D \ten s_D'), \\
      - D &\coloneqq (L_D^{-1}, |D|, s_D^{-1}).
    \end{align*}

    For a morphism $f:Y \to X$, define $f^* D \coloneqq (f^* L_D, f^{-1} (|D|), f^* s_D)$.
  \end{definition}

  \begin{definition}
    When $X$ is an $n$-dimensional variety, define $[D] \in \CH_{n-1} (|D|)$ to be $c_1(L_D)|_{|D|}$.
  \end{definition}

  \begin{definition}
    Let $X$ be a scheme, $D$ be a pseudo-divisor on $X$, $V \subset X$ be a subvariety of $X$ of dimension $i$ with inclusion $j: V \to X$.
    We define 
    \[ D \cdot [V] \coloneqq [j^* D] \in \CH_{i-1} (|D| \cap V). \]

    For an $i$-cycle $\alpha = \sum n_V [V]$ on $X$, define 
    \[ D \cdot \alpha \coloneqq \sum n_V (D \cdot [V]) \in \CH_{i-1}(|D| \cap |\alpha|), \]
    where $|\alpha| \coloneqq \bigcup_{n_V \neq 0} V$ 
    is the support of $\alpha$.
  \end{definition}

  \begin{example}
    Let $X$ be a purely $n$-dimensional scheme, $D$ be a Cartier divisor on $X$.
    Then $D \cdot [X] = [D]$ by \cref{lem:restriction_of_cartier_divisor_to_irreducible_components}.
  \end{example}

  \begin{proposition}
    We have the following basic properties.
    \begin{enumerate}
      \item Let $X$ be a scheme, $D,D'$ be pseudo-divisors on $X$, $\alpha,\alpha'$ be $i$-cycles on $X$.
        Then 
        \[ D \cdot (\alpha + \alpha') = D \cdot \alpha + D \cdot \alpha' \in \CH_{i-1}(|D| \cap (|\alpha| \cup |\alpha'|)). \]

      \item The same as (a). 
        We have 
        \[ (D + D') \cdot \alpha = D \cdot \alpha + D' \cdot \alpha \in \CH_{i-1}\big((|D| \cup |D'|) \cap |\alpha|\big). \]

      \item Let $f: X \to Y$ be a proper morphism, $\alpha$ be an $i$-cycle on $X$, $D$ be a pseudo-divisor on $Y$, 
        then $f_* (f^* D \cdot \alpha) = D \cdot f_* \alpha \in \CH_{i-1}(|D| \cap f(|\alpha|))$.

      \item Let $f:X \to Y$ be a flat morphism of relative dimension $n$, $\alpha$ be an $i$-cycle on $Y$, $D$ be a pseudo-divisor on $Y$, then 
        \[ f^* (D \cdot \alpha) = f^* D \cdot f^* \alpha \in \CH_{i+n-1}(f^{-1} (|D|) \cap f^{-1}(|\alpha|)). \]

      \item Let $D$ be a pseudo-divisor on a scheme $X$ s.t. $L_D \cong \calO_X$.
        Then, for all $i$-cycles $\alpha$ on $X$, $D \cdot \alpha = 0 \in \CH_{i-1} (|\alpha|).$
    \end{enumerate} 
  \end{proposition}
  \begin{proof}
    (a) By definition.

    Hence in (b)-(e), may assume $\alpha = [V]$ for some variety $V$.

    (b) By linearity of restricting pseudo-divisors to a subvariety and taking the class of pseudo-divisors.

    (c) If $f$ is a closed embedding, it is ok.
    In general, may assume: $X$ is a variety, $\alpha = [X]$, $f:X \to Y$ is surjective proper, $D$ is a Cartier divisor.
    WTS, $f_* f^* [D] = \deg (X/Y) [D]$.
    By working locally, may assume $D = \divisor(r)$ for some $r \in \fkk (X)^\times$, then 
    \[ f_* f^* \divisor(r) = \divisor(r^d) = d \divisor(r), \]
    where $d = \deg (X/Y)$.

    (d) May assume: $Y$ is a variety and $\alpha = [Y]$, 
    WTS: $f^* [D] = [f^{-1}(D)]$.
    It is by \cref{lem:flat_pullback_of_scheme_is_inverse_image}.

    (e) May assume $X$ is a variety, $\alpha = [X]$, $D = \divisor(r)$ for some $r \in \fkk(X)^\times$, then $\divisor(r) \cdot [X] = \divisor(r) = 0 \in \CH_{i-1}(X)$.
  \end{proof}

  \begin{theorem}
    Let $D,D'$ be Cartier divisors on an $n$-dimensional variety $X$.
    Then 
    \[  D \cdot [D'] = D' \cdot [D] \in \CH_{n-2}(|D| \cap |D'|). \]
  \end{theorem}
  \begin{proof}
    \hspace{1em}
    \begin{case}
      $D,D'$ are effective Cartier divisors (i.e. codimension $1$ closed subschemes, locally defined by single non zero functions).
    \end{case}
    By induction on:
    \[ \calE(D,D') \coloneqq \max \{ \ord_V (D) \ord_V (D') \mid V \subset X \text{ of codimension } 1 \}. \]
    The basis case is $\calE(D,D') = 0$, i.e. $D$ and $D'$ intersect properly.
    Let $W \subset X$ a codimension $2$ subvariety, $A \coloneqq \calO_{X,W}$, $a, a' \in A$ be local equations of $D$ and $D'$.
    Let $V \subset X$ be a codimension $1$ subvariety containing $W$ which correspond to a height $1$ prime ideal $\frakp \subset A$.
    Then the 
    \begin{itemize}
      \item coefficient of $[V]$ in $[D']$ is $\call(O_{D',V}) = \call_{A_\frakp} (A_\frakp / a' A_\frakp)$, and
      \item coefficient of $[W]$ in $D \cdot [V]$ is $\call(O_{D \cap V,W}) = \call_{A/\frakp} (A/(\frakp + a A))$.
    \end{itemize}
    Hence the coefficient of $[W]$ in $D \cdot [D']$ is 
    \[ \sum_\frakp \call_{A_\frakp}(A_\frakp / a' A_\frakp) \call_{A/\frakp}(A/(\frakp + a A)). \]
    By \cite[A 2.7 and A 2.8]{Fulton13Intersection}, change $a$ and $a'$ is ok.
    By symmetry, $D \cdot [D'] = [D'] \cdot [D]$ as cycles.

    Now we assume $\calE(D,D') > 0$ (just an idea).
    Let $\pi: \tilde{X} \to X$ be the blow-up along $D \cap D'$, $E \subset \tilde{X}$ be the exceptional divisor.
    Then $\pi^* D = \tilde{D} + E, \pi^* D' = \tilde{D'} + E$, where $\tilde{D}$ and $\tilde{D'}$ are effective Cartier divisors such that $\tilde{D} \cap \tilde{D'} = \emptyset$.
    Moreover, $\calE(\tilde{D}, E), \calE(\tilde{D'}, E) < \calE(D, D')$ (for the proof, read Fulton).
    Then 
    \[ D \cdot [D'] = \pi_*((\tilde{D} + E) \cdot ([\tilde{D'}] + [E])) = \pi_*((\tilde{D'} + E) \cdot ([\tilde{D}] + [E])) = D' \cdot [D]. \]

    \begin{case}
    $D$ is arbitrary, $D'$ is effective.
    \end{case}

    If $D$ is the difference of $2$ effective Cartier divisors, the proof is reduced to case 1.
    However, this is in general FALSE.

    In general, may assume: $X = \Spec A$, $a \in \fkk(X)^*$ be the local equation for $D$.
    Define $I \coloneqq \{b \in A \mid a \cdot b \in A\}$ be the ideal of denominations of $a$.
    Let $\pi: \tilde{X} \to X$ be the blow-up along $V(I)$, $E$ be the exceptional divisor.
    Then $\pi^* D = C - E$, where $C$ is an effective Cartier divisor.
    Since the Theorem holds for $(C , \pi^* D'), (E, \pi^* D')$, it also holds for $(D,D')$.

    \begin{case}
      $D,D'$ are arbitrary.
    \end{case}

    Let $\pi:\tilde{X} \to X$ as Case 2.
    Since the Theorem holds for $(C, \pi^*D')$ and $(E,\pi^*D')$.
    It also holds for $(D,D')$.
  \end{proof}

  \Yang{I need more computation for the theorem, in particular the blow up trick.}

  \begin{corollary}\label{cor:intersection_with_rationally_zero_cycle_is_zero}
    Let $X$ be a scheme, $\alpha$ be an $i$-cycle on $X$.
    If $\alpha \sim 0$ on $X$, then for all pseudo-divisor $D$ on X, $D \cdot \alpha = 0 \in \CH_{i-1}(|D| \cap |\alpha|)$.
    In particular, we have well defined 
    \[ \CH_i (X) \to \CH_{i-1} (|D|),\quad \alpha \mapsto D \cdot \alpha. \]
  \end{corollary}
  \begin{proof}
    May assume: $X$ is a variety, $\alpha = \divisor(f)$ for some $f \in \fkk(X)^*$. 
    Then 
    $ D \cdot \divisor(f) = \divisor(f) \cdot [D] = 0 \in \CH_{n-2}(|D| \cap |\alpha|). $
  \end{proof}

  \begin{definition}
    Let $X$ be a scheme, $L$ be a line bundle on $X$.
    For an $i$-dimensional subvariety $V$ of $X$, define 
    \[ c_1(L) \cap [V] \coloneqq \text{the class of } L|_V \in \CH_{i-1}(V). \]
    By \cref{cor:intersection_with_rationally_zero_cycle_is_zero}, we have well defined 
    \[ c_1(L)\ \cap : \CH_i(X) \to \CH_{i-1} (X). \]
  \end{definition}